% !TITLE 江南逢李龟年
% !AUTHOR 杜甫
% !DYNASTY 唐
% !ORDER 31

\begin{poem}
岐王宅里寻常见\textsuperscript{1}，崔九堂前几度闻\textsuperscript{2}。
正是江南好风景，落花时节又逢君\textsuperscript{3}。
\end{poem}

\begin{annotations}
\notetext{1}{岐王宅：岐王李范的宅第。李范是唐玄宗的弟弟，喜好文艺，常在家中宴请文士。“寻常见”——在岐王府上经常见到你，当年的相见是那样平常，不值得珍惜。}
\notetext{2}{崔九堂：殿中监崔涤（排行第九）的厅堂。李龟年常在崔涤府上演唱。“几度闻”——在崔九家中多次听到你的歌声。}
\notetext{3}{落花时节又逢君：在落花缤纷的暮春，又遇见了你。四十年前在长安盛世相逢，四十年后在江南落花中重逢——“又”字里有千钧之重。}
\end{annotations}

\begin{appreciation}
这是杜甫最后一首名作，写于大历五年（770年），就在他去世的那一年。在潭州（今长沙），他偶然遇到了开元时代的歌唱家李龟年。

李龟年是盛唐最著名的音乐家，深受唐玄宗赏识，经常在岐王李范和殿中监崔涤的府第演唱。杜甫年轻时曾在这些贵族的宴会上见过他，听过他的歌声。那是开元盛世，长安繁华，歌舞升平。

岐王宅里寻常见，崔九堂前几度闻——在岐王的宅子里常常看见你，在崔九的厅堂前多次听到你的歌声。寻常和几度都是轻描淡写的词，但背后是一个时代的记忆。那时的见面是轻松的、频繁的，不需要珍惜，因为以为永远都会这样下去。

正是江南好风景，落花时节又逢君——现在正是江南风景好的时候，在落花缤纷的时节，又遇到了你。这两句看似平淡，实则千钧之重。江南好风景——风景确实好，但人已经老了，国已经破了。落花时节——春天即将过去，花朵正在凋零。这不仅是写景，是写时代：盛唐已经过去了，像落花一样不可挽回。

又逢君——一个又字，藏着多少感慨。四十年前在长安相逢，那时两人都年轻，都处在最好的时光里；四十年后又在江南相逢，两人都已垂老，盛世已经成了一场梦。杜甫没有写自己的悲伤，他只是说落花时节又逢君。但读者都能感受到：在那个落花时节，两个老人相对无言，心里想的都是同一个问题——那些好日子，都到哪里去了？

这首诗是中国文学史上最克制的抒情诗之一。二十八个字，没有一句直接写悲，却字字都是悲。这就是杜甫晚年达到的境界：最高的技巧，是不见技巧；最深的情感，是不说情感。
\end{appreciation}
